\section{Limitations}
\label{sec:limitations}

The empirical scope is one simulated greenhouse-tomato domain. The mechanistic crop model
improves traceability, but the economic coefficients, site distribution, factor ranges,
facility layout and observation channels are benchmark choices. Results do not establish
the same tool effects in chemistry, materials, biology, field agriculture or a physical
greenhouse.

A preregistered observed-action GreenLight v8 holdout also failed its trajectory gate:
pooled temperature RMSE was 2.579\,$^\circ$C against a 2.04\,$^\circ$C limit and
relative-humidity RMSE was 9.983 percentage points against an 8.5-point limit; CO$_2$
passed its separate threshold. We therefore use public greenhouse trajectories to bound
observed ranges, dynamics, missingness and executor sensitivities, not to claim a validated
digital twin or counterfactual actuator validity. Realised fogging flow, heating-water
mass flow, natural-ventilation airflow and screen heat/moisture flux remain unobserved in
the audited public records.

The prospective model extension covers Luna, DeepSeek V4.1 Flash and Qwen 3.8 27B, one
prompt family and one inference-aid implementation. Provider-hidden reasoning was disabled
for all three. DeepSeek used its direct API with native JSON mode, whereas Luna and Qwen
used OpenRouter; Qwen's requested and returned model identifier was stable, but OpenRouter
routed calls across multiple underlying providers. The effects therefore describe these
tested model--interface--provider conditions, rather than provider-invariant model
properties. The earlier Sol replication uses independent sites and remains supplementary;
it cannot be included in the same-site model ranking.

The 288-episode model comparison was prospectively frozen only after the original Phase~4
analysis. Interface-only DeepSeek pilots were excluded before inspecting effects, and the
formal amendment and analysis were frozen before the extension was analyzed. This protects
the extension from outcome-driven choices but does not make it part of the original
confirmatory study. Replication on newly reserved sites is needed, especially for the Qwen
interaction.

Within-cycle observations are consequential only through the recommendation during the
current synchronous protocol. Units cannot be released, treatments cannot be modified, and
the next batch begins after terminal outcomes are available. The experiment therefore tests
whether inspection changes reasoning and endpoint recommendations, not the value of early
stopping, asynchronous replacement or adaptive treatment. Those actions require explicit
cleanup time, residual value and staggered-capacity semantics.

The common-history evaluator has a strong one-factor calibration but incomplete
higher-dimensional calibration. Its finite candidate search also supplies a feasible lower
bound on the world optimum. We consequently base primary comparisons on realized regret and
restrict posterior-risk decompositions to effects larger than measured evaluator error.

The experimental \textsc{Screen} configuration is not yet a screening benchmark. It asks
for an operating recommendation and scores final regret, without eliciting factor rankings,
effect directions, calibrated confidence or top-$k$ selections. We therefore exclude it
from the core suite and defer screening claims until those outputs and direct metrics exist.

Finally, the Phase~3 scripted reference rows and Phase~4 LLM rows use independent held-out
site ranges. They provide common-scale descriptive anchors, but they are not paired
model-versus-baseline tests. Direct superiority claims require a jointly randomized matrix
on the same sites.
